\documentclass[../main]{subfiles}
\begin{document}
\ifSubfilesClassLoaded{\setcounter{page}{14}}{}
\section{Основание коммунизма}
\label{sec:03_osnovanie}

Основанием для капитализма является товарное производство. Всё богатство капиталистического общества выступает как совокупность товаров. Поскольку продукты труда при капитализме неизбежно приобретают товарную форму, товар является «клеточкой капитализма».

Основанием для коммунизма как общественного устройства, в соответствии с материалистическим пониманием истории, также может быть только соответствующий способ производства.

Всё богатство коммунистического общества - универсальность содержания и форм человеческой деятельности. Марксизм вообще видит сущность человека в деятельности по изменению природы (в том числе и своей собственной общественной природы). Трудиться, обрабатывать природный материал, создавать вещи, которых нет в природе - вот что принципиально отличает жизнь человека от жизнедеятельности животного и определяет все остальные отличия. Но до сих пор сам труд в общественном масштабе был несвободным. Все докоммунистические способы производства отличались своей специфической формой несвободы труда. Коммунизм же предполагает полное преодоление этой несвободы. Поэтому его основанием может быть только свободный труд «без нормы и вознаграждения», как «первейшая жизненная потребность» человека.

Простейшая непосредственно-коллективная форма осуществления такого труда (а лучше сказать, такой деятельности, так как за словом «труд» в известной мере сохраняется негативное значение) является той самой «клеточкой» коммунизма, в которой заложено дальнейшее развитие коммунистического общества во всех своих необходимых определениях.

Свободная деятельность и её непосредственно-коллективная форма являются реальностью в наше время точно так же, как реальностью ещё до появления капитализма были товарное производство и товар. И переход к новому общественному строю произойдет, когда свободная деятельность человека станет универсальным отношением, так же, как таковым некогда стал товар. Из пестроты имеющихся на рынке товаров, обладающих различными свойствами, нельзя непосредственно понять, что такое товар и товарное производство. Из разнообразия видов и способов организации коммунистической деятельности в их особенности нельзя понять, что она такое и что такое её непосредственно- коллективная форма. Тем более, что слова могут заслонить суть дела. Непосредственно-коллективная форма свободной деятельности может иметь место даже тогда, когда человек осуществляет её совершенно один. Нужно выделить те определения, которые являются необходимыми для неё.

В    непосредственно-коллективной      форме      свободная    деятельность осуществляется по меркам свободы, истины, добра и красоты, а значит и наивысшего человеческого интереса, то есть - это форма самоценной чувственной деятельности. Свободная деятельность обязательно предполагает единство всех трёх компонентов: чувственность, нравственность и мышление, и доведение их до уровня свободы. В процессе такой деятельности осуществляется утверждение человека во всей его полноте вне зависимости от того, что именно делается и что именно (истина, добро или красота) выходит на передний план.

Такая непосредственно-коллективная форма - это всего лишь взятая со стороны формы сама эта деятельность с соответствующим ей содержанием. В обычной речи такую деятельность часто называют социальным творчеством или просто творчеством. Форма является важной, так как развитие осуществляется и выражается в необходимом движении формы. Содержание всегда вызывает к жизни соответствующую форму и осуществляется через неё. Поэтому противоречия содержания схватываются сознанием первоначально как противоречивость формы. В её развёртывании и через её развертывание (смену форм) в действительности происходит движение содержания. Со стороны формы в деятельности важны такие моменты, как характер связи между людьми, цель, общественное бытие предмета и средств труда.

Непосредственная коллективность в свободной деятельности выражается, в первую очередь, в нацеленности на других людей. В непосредственно- коллективных формах деятельности человек выступает всегда как цель и наивысшая ценность. Всё, что делается, делается непосредственно для развития других людей. Поэтому и продукт такого труда никогда не приобретает самостоятельно-вещную форму (форму товара). Продукт такого труда не может противостоять живым людям как чужая им вещь. Другими словами, эта деятельность направлена на развитие общества и улучшение жизни других людей, осуществляется для них, а потому и для себя, как представителя этого общества. То, что сейчас принято называть материальными ценностями в свободной коммунистической деятельности может служить лишь средством.

Что же касается характера связи между людьми, то её в общем и целом можно определить как товарищество, или, что то же самое, свободная ассоциация.

Этот характер связи указывает также и на средства такой деятельности, так как в свободной ассоциации ни один из участников не может быть средством для осуществления чужой воли и чуждых ему интересов. Даже если кто-то, благодаря своим способностям и умениям, является ведущим в коллективе, остальные не слепо исполняют чужую волю, а сотрудничают в осуществлении общего дела, одинаково важного и нужного для всех его участников.

Характер связи также указывает на общественное бытие средств и предмета коммунистической деятельности. Таким образом, в качестве средства могут выступать только отличные от человека вещи и процессы. Общественная форма бытия средств и предметов труда предполагает их непосредственно- общественное использование. Именно в этом смысле Маркс говорил о переходе от управления людьми к управлению вещами.

Из сказанного выше видно, что такие формы человеческой деятельности лишены какого бы то ни было внешнего как прямого, так и непрямого (вознаграждение за труд) принуждения, поскольку ни процесс, ни результат коммунистической деятельности не является безразличным для того, кто её осуществляет. Человек всякий раз в своём труде заинтересован непосредственно, в отличии от нормальной для современного общества ситуации, когда сама деятельность выступает лишь средством для получения денег.

Конечно же, непосредственная коллективность заключается ещё и в том, что труд любого члена общества опирается на труд других людей, в том, что предмет, средства, и способ труда является продуктом, выработанным другими людьми, а его собственный труд, в свою очередь, является неотъемлемой частью всего общественного труда и даёт материал для будущей деятельности. Но это характерно для любого труда, осуществляемого в любой общественной форме, а не только для коммунистической деятельности, а потому не касается специфики коммунизма.

Под непосредственно-коллективной формой свободной деятельности мы понимаем объединение людей в процессе деятельности, где целью каждого является развитие всех, а средствами для её достижения выступают вещи и природные процессы, по отношению к которым происходит свободное и равное волеизъявление.

Непосредственно-коллективная форма деятельности — это не просто внешняя общественная форма деятельности, которая может иметь место, а может и не иметь, не затрагивая при этом содержания деятельности, а напротив, - это её содержательная форма. Вне её вообще невозможно осуществление свободной деятельности по меркам истины, добра и красоты.\footnote{Имеется в виду вовсе не расхожее представление, где красота связывается со вкусом, добро с мнениями о том, что такое хорошо и плохо, которые могут быть разными, и только истине придается какое-то объективное значение. Категории <<истины>>, <<добра>> и <<красоты>> разрабатывались на протяжении тысячелетий, как объективные идеальные характеристики человеческой деятельности в мире, разворачивающиеся (развивающиеся) в этой деятельности, а потому и объективные характеристики мира (материи). Другими словами, если коротко, когда говорится о деятельности по меркам истины добра и красоты, то речь идёт о <<точках развития>> человеческой чувственности, познания и нравственности, о развитии отношения человека к человеку и человека к природе, взятых со стороны идеального измерения.} Движение и метаморфозы такой формы предполагают сохранение вышеозначенных моментов, подобно тому, как в метаморфозах товарной формы продуктов труда всегда сохраняется стоимость (то, что делает товар именно товаром, а не просто продуктом труда).

Что же касается содержания такой деятельности, она - всегда революционна. В этом заключается её противоречие по отношению к самой себе. В свободной коммунистической деятельности постоянно полагается и её внутренний предел, и выход за этот предел. С помощью исторически достигнутых, а значит исторически ограниченных рубежей деятельности, меняется сама эта деятельность, полагается её новое содержание (благодаря расширению горизонта освоения материи) и новая форма. Это происходит не независимым от участников процесса труда окольным путём, как это происходит в процессе вещного производства. Наоборот, это изменение происходит целенаправленно. Поэтому коммунизм, как освоение человеком своей собственной    деятельности,    обнажает    внутреннюю     противоречивость человеческой деятельности вообще, а значит и внутреннюю противоречивость материального мира, в котором человек осуществляет свою деятельность. Потому и для науки о коммунизме важнейшее методологическое значение играет замечание Маркса, высказанное в первом тезисе о Фейербахе. «Главный недостаток всего предшествующего материализма - включая и фейербаховский - заключается в том, что предмет, действительность, чувственность берется только в форме объекта, или в форме созерцания, а не как человеческая чувственная деятельность, практика, не субъективно. Отсюда и произошло, что деятельная сторона, в противоположность материализму, развивалась идеализмом, но только абстрактно, так как идеализм, конечно, не знает действительной, чувственной деятельности как таковой. Фейербах хочет иметь дело с чувственными объектами, действительно отличными от мысленных что такое хорошо и плохо, которые могут быть разными, и только истине придается какое-то объективное значение. Категории “истины”, “добра” и “красоты” разрабатывались на протяжении тысячелетий, как объективные идеальные характеристики человеческой деятельности в мире, разворачивающиеся (развивающиеся) в этой деятельности, а потому и объективные характеристики мира (материи). Другими словами, если коротко, когда говорится о деятельности по меркам истины добра и красоты, то речь идёт о “точках развития” человеческой чувственности, познания и нравственности, о развитии отношения человека к человеку и человека к природе, взятых со стороны идеального измерения. объектов, но самое человеческую деятельность он берет не как предметную деятельность». В этом замечании не только подчеркнута роль идеализма в становлении теории деятельности. Хотя эта роль действительно огромна и важна, настолько, что впоследствии Ленин говорил о важности развития обществ «материалистических друзей гегелевской диалектики». Материалистически понятые категории диалектики как «миллионы раз повторенные формы человеческой деятельности» (Ленин), их внутренняя противоречивость и движение являются необходимым инструментом развития теории коммунизма именно потому, что в них отражены существенные моменты человеческой деятельности вообще, схвачено то, без чего человеческая деятельность в принципе не может осуществляться. Но в Марксовом тезисе также делается акцент на материальности человеческой деятельности.

Подчеркнем ещё раз, что свободная деятельность - это не «отдых» в современном смысле слова, не праздное провождение времени, а деятельность по меркам свободы. Когда-то философия в лице Бенедикта Спинозы определила свободу как осознанную необходимость. Настоящая свобода определяется действием сообразно с этой необходимостью.

В свободное время в соответствии c исторической необходимостью люди развивают культуру в самом широком смысле этого слова. Такая деятельность не является придуманной коммунистами в фантастических мечтах. Она существовала в прошлом и существует сейчас. Но она существовала и существует в очень жёстких рамках и является достоянием лишь немногих людей. Абсолютное большинство членов общества не знает свободной деятельности. Это происходило и происходит, в том числе, потому, что нужно не просто иметь время, но и быть развитым настолько, чтобы быть в состоянии такую деятельность осуществлять. Условия для такого развития были и остаются в современном обществе товарного производства весьма ограниченными.

Итак, чтобы двигаться в направлении коммунизма, производить нужно коммунизм, то есть свободную деятельность, точнее, свободное время и соответствующий уровень развития членов общества, которые могли бы её осуществлять.

Следовательно, и распределение здесь, в первую очередь, это распределение свободного времени и возможностей освоения человеческой культуры. Распределение продуктов труда необходимо лишь постольку, поскольку оно обеспечивает свободное время. В этом духе должен ставиться вопрос об эффективности общественного производства в переходный период к коммунизму. Свободное время при капитализме является издержками производства, оно не учитывается, не ценится, и потому часто так и не становится действительно свободным, не используется по назначению. Назначение свободного времени - быть использованным для постоянного развития общественных отношений (или развития деятельности, что, в сущности, одно и то же).

Потому для коммунизма привычка, ставшая достоянием всех, относиться к собственному свободному времени и деятельности, которая осуществляется в это время, как к общественному богатству, не подлежащему частному разбазариванию, является очень важной. Вообще, сила привычки и коллективного формирования привычки имеет здесь огромное значение, хотя бы потому, что быть свободным, то есть действовать по меркам свободы, всегда очень трудно. Это означает находиться на острие противоречий, через которые осуществляется развитие общества. Это не только тяжёлый труд, но и трагедия в том смысле, что нужно постоянно преодолевать себя, отрицать себя прошлого. А коллективная общественная свобода - это коллективная общественная трагедия и её постоянное разрешение. Коммунистическое общество - это не общество всеобщей лени и праздности, а общество всеобщей свободной деятельности, являющейся для людей первой жизненной потребностью.

Коммунизм осуществляется через необходимую смену непосредственно коллективных форм человеческой деятельности по изменению самой этой деятельности (развитию самой жизни людей). Достигнутые формы, стремясь за свои пределы, создают новые. Потому эту смену нельзя представить как постоянно повторяющиеся следующие друг за другом одни и те же метаморфозы (как это мы видим в движении капитала). Внутренняя логика движения непосредственно-коллективных форм свободной деятельности, необходимость их смены определяется не просто исторически достигнутым способом производства и соответствующим им производственным отношениями, а свободой как необходимостью самой материи, раскрывающейся по мере освоения материи.

В наиболее общем виде закон коммунизма можно сформулировать как закон свободы или закон самодетерминации человеческой деятельности: свободная деятельность членов общества, осуществляющаяся как сознательное развитие содержания и форм этой деятельности, определяет развитие всего общества. Или, другими словами, свобода, реализуясь, стремится к самовозрастанию.

Конечно же, понимание действия этого закона требует постоянного уточнения с использованием всех имеющихся знаний, как о природе, так и об обществе. Потому теория коммунизма может в соответствующих условиях играть объединяющую роль для всех наук, указывая на реальную связь между ними, как связь самой человеческой практики.

Существующие в наши дни непосредственно-коллективные формы свободной человеческой деятельности как «ростки коммунизма» не просто являются свидетельством его возможности, они есть его реальное бытие. Важно понять, что коммунистический способ производства уже есть, речь идет лишь о том, чтобы он стал господствующим и распространился на всё общество.

В своём движении непосредственно-коллективные формы свободной человеческой деятельности как «ростки коммунизма» могут развернуться в коммунизм как общественный уклад, то есть как способ существования общества. В этом смысле одним из наиболее ценных положений в теоретическом наследии Ленина является совет выявлять и лелеять «ростки» коммунизма везде, где они есть.
\end{document}
